\section{Related work}
\label{sec:study}

\paragraph{World models and inverse dynamics.}
UniSim~\cite{yang2023unisim} learns a simulator of real-world interaction from mixed data and trains policies inside it.
IRASim~\cite{zhu2024irasim} generates robot videos from a first frame and end-effector actions, conditioning each frame on its own action, and is used for policy evaluation with human-judged rollouts and for planning.
Cosmos~\cite{nvidia2025cosmos} provides video world foundation models, one of them post-trained on Bridge with actions.
Inverse dynamics models, which predict the action between two frames, label unlabelled videos in VPT~\cite{baker2022vpt}, turn video plans into robot actions in UniPi~\cite{du2023unipi}, and label generated videos as training data in DreamGen~\cite{jang2025dreamgen}; these works trust the actions read, while we compare them with the commanded ones.

\paragraph{Evaluating generated videos.}
VideoPhy~\cite{bansal2024videophy} and WorldModelBench~\cite{li2025worldmodelbench} judge physical commonsense and instruction following; Physics-IQ~\cite{motamed2025physicsiq} compares predicted continuations of real experiments with the recorded ones and finds that realism does not imply physical understanding; EWMBench~\cite{yue2025ewmbench} compares end-effector trajectories with the ground truth and notes that static videos score high on scene consistency.
WorldSimBench~\cite{qin2024worldsimbench}, World-in-World~\cite{zhang2026worldinworld}, and WoW-World-Eval~\cite{fan2026wowwoval} read actions from generated videos and execute them in simulators or on robots, scoring task success; KineBench~\cite{liu2026kinebench} replaces the inverse dynamics model by pose estimation because its errors confound the result; SIMPLER~\cite{li2024simpler} matches simulators to real setups, including Bridge.
\ours{} needs neither the real video nor a simulator, and it measures the error of its reader to separate it from that of the world model.
\Cref{tab:matrix} in the appendix compares the main sources.

\paragraph{Open questions.}
Does a realistic frame imply that the robot followed the commanded actions, do objects move only when touched and keep their identity, and could the real arm execute the motion?
We measure the first and the third with \ours{}, and the second only through vision-language judges.

% ---- review comments, kept for the record (text above rewritten for concision on 2026-11-02) ----
% (Turn 2) "where humans judge the success of generated rollouts": This part want to show the idea where IRASim generate the videos from the action trajactories generated by the policy, then use human to judge these generated videos by IRASim to decide which on is the best, right? If so, I think we should use some easier-to-understand words to express this
% Claude: partly. The first half is right: a policy produces the actions, IRASim generates the videos, humans judge them. But the humans do not pick the best trajectory: they label every generated rollout success or failure, and the success rate in IRASim is compared with the true simulator's (Pearson 0.99 over four policy checkpoints). Picking the best is a second experiment, planning, where a learned reward model scores the generated outcomes. Both are now explained in plain words, in their own sentence.
% (Turn 2) "An inverse dynamics model predicts the action between consecutive observation": Don't you think this claim which stands alone is out of place/context?
% Claude: agreed; it now opens its own paragraph and says why it matters (it is how videos become actions), and the paragraph ends with how our use differs.
% (Turn 2) " VPT [1] trains one on a small labelled set to label large unlabelled video collections, UniPi [3] uses one to turn text-conditioned video plans into robot actions, and DreamGen [5] uses one to label videos generated by a world model as training data for visuomotor policies.": You repeated word "one" here, which looks a little bit weird, in professional paper they do not write like this I think since the reader will be confused what is "one"
% Claude: agreed; "such a model" once, then "it", after the paragraph has named inverse dynamics models.
% (Turn 2) "Physics-IQ": What does this work do? Has it already solved our raised problems?
% Claude: it tests whether a video model understands passive physics (solids, fluids, optics) by continuing real experiment videos and measuring motion overlap with the true continuation. It does not solve our problem: it needs the real continuation, has no commanded actions, and no robot. It supports our first observation (looking real is not being physically right). Now said in the text.
% (Turn 2) "SIMPLER [7] evaluates real-robot policies in simulated scenes matched to real setups, including the Bridge setup.": How does this work relate to our homework?
% Claude: two links, now written out: it is the physics-simulator alternative that world models want to replace (the reason the brief asks whether generated motion is executable), and it simulates the same WidowX robot as our clips, so it is the natural next step after our kinematic executability check (replay with dynamics and contacts).
% Claude: follows from the executability check now being run (was "propose a test for the third").
